\Askhsh{24587}{4}
\begin{ylh}{1}
    \item 2.1 Η έννοια της παραγώγου
    \item 2.2 Παραγωγίσιμες συναρτήσεις - Παράγωγος συνάρτηση
    \item 2.3 Κανόνες παραγώγισης
    \item 2.6 Συνέπειες του Θεωρήματος Μέσης Τιμής
    \item 2.7 Τοπικά ακρότατα συνάρτησης.
\end{ylh}
Δίνεται η συνάρτηση $f:(0,+\infty) \rightarrow \mathbb{R}$, με τύπο $f(x) = 2 \ln x - x$ και η ευθεία $\varepsilon: y = x$.
Γνωρίζουμε ότι η απόσταση ενός σημείου $M(x_0,y_0)$ της γραφικής παράστασης της συνάρτησης $f$ από την ευθεία $\varepsilon$, είναι $d(M,\varepsilon) = \sqrt{2} |x_0 - \ln x_0|$.
\begin{alist}
\item Να αποδείξετε ότι η απόσταση του σημείου $M(x_0,y_0)$ της γραφικής παράστασης της συνάρτησης $f$ από την ευθεία $\varepsilon: y = x$, είναι $d(M,\varepsilon) = \sqrt{2} (x_0 - \ln x_0)$. \monades{05}
\item \begin{rlist}
    \item Να βρείτε το σημείο της $C_f$, το οποίο απέχει την ελάχιστη απόσταση από την ευθεία $\varepsilon$. \monades{12}
    \item Να βρείτε την ελάχιστη απόσταση. \monades{03}
\end{rlist}
\item Να βρείτε το σημείο της $C_f$ στο οποίο η εφαπτομένη της είναι παράλληλη με την ευθεία $y = x$ και στη συνέχεια να βρείτε την εξίσωση της εφαπτομένης. \monades{05}
\end{alist}